\section{Prezentowalne \texorpdfstring{$\infty$-kategorie}{kategorie wyższe} i lokalizacje Bousfielda}\label{cat:sec:presentable-infty-bousfield}

W rozdziale~\ref{cat:sec:accessible-presentable} badaliśmy klasyczne kategorie dostępne i lokalnie prezentowalne: ich obiekty można budować z małej rodziny obiektów za pomocą odpowiednich kogranic filtrowanych. W rozdziale~\ref{cat:sec:derived-localization-morita} odwracaliśmy morfizmy w $\infty$-kategoriach. Połączymy teraz te idee. Chcemy rozumieć, kiedy duża $\infty$-kategoria jest generowana przez mały zasób obiektów i kiedy lokalizacja ma pełny model w postaci podkategorii obiektów lokalnych. W odróżnieniu od klasycznej teorii kontrolujemy nie tylko klasy morfizmów, lecz całe przestrzenie odwzorowań.

\subsection{Filtrowanie wyższych kategorii}

Ustalamy uniwersum, względem którego rozumiemy małość, i nieskończoną regularną liczbę kardynalną $\kappa$ (por. definicja~\ref{cat:def:accessible-regular}). Przez $\mathcal S$ oznaczamy $\infty$-kategorię przestrzeni, modelowaną przez kompleksy Kana z odwróconymi słabymi równoważnościami homotopijnymi.

\begin{definicja}[$\kappa$-filtrowana $\infty$-kategoria]\label{cat:def:infty-accessible-filtered}
Zbiór symplicjalny $K$ nazywamy $\kappa$-małym, jeśli ma mniej niż $\kappa$ sympleksów niezdegenerowanych. $\infty$-kategorię $\mathcal I$ nazywamy $\kappa$-filtrowaną, jeżeli każde odwzorowanie $K\to\mathcal I$ z $\kappa$-małego zbioru symplicjalnego przedłuża się do $K^{\triangleright}\to\mathcal I$, gdzie $K^{\triangleright}=K\star\Delta^0$ jest stożkiem górnym. Dla $\kappa=\aleph_0$ mówimy po prostu o filtrowaniu.
\end{definicja}

Diagram pusty również jest dopuszczalny, więc $\mathcal I$ jest niepusta. Wybór stożka obejmuje nie tylko morfizmy do wspólnego obiektu, ale też homotopie między złożeniami oraz zgodności wyższych rzędów. W szczególności nie wystarcza żądać istnienia wspólnego celu dla skończenie wielu obiektów.

\begin{stwierdzenie}[Porównanie z kategoriami zwykłymi]\label{cat:prop:infty-filtered-ordinary}
Dla zwykłej małej kategorii $I$ nerw $NI$ jest $\kappa$-filtrowany wtedy i tylko wtedy, gdy $I$ jest $\kappa$-filtrowana w sensie definicji~\ref{cat:def:accessible-filtered}. Jeżeli $\mathcal I$ jest $\kappa$-filtrowaną $\infty$-kategorią, to $h\mathcal I$ jest $\kappa$-filtrowaną zwykłą kategorią.
\end{stwierdzenie}
\begin{proof}
Załóżmy najpierw, że $I$ jest $\kappa$-filtrowana. Dany diagram $K\to NI$ wyznacza rodzinę obiektów $X_v$ indeksowaną wierzchołkami i morfizmów $f_e:X_v\to X_w$ indeksowaną niezdegenerowanymi krawędziami. Obie rodziny mają moc mniejszą niż $\kappa$. Wybieramy wspólny cel $Y$ i morfizmy $a_v:X_v\to Y$. Dla każdej krawędzi $e:v\to w$ mamy parę równoległych morfizmów $a_v,a_wf_e:X_v\to Y$. Możemy jednocześnie wyrównać wszystkie pary jednym morfizmem $b:Y\to Z$: każda para ma mniej niż $\kappa$ elementów, a rodzina par ma moc mniejszą niż $\kappa$, więc stosujemy kryterium filtrowania z twierdzenia~\ref{cat:prop:accessible-filtered-criterion}. Wtedy $ba_v=ba_wf_e$. Dla każdego sympleksu $K$ uzupełniamy odpowiadający mu ciąg składalnych strzałek w $I$ o strzałkę z ostatniego obiektu do $Z$. Równości dla krawędzi zapewniają zgodność ze ścianą usuwającą ostatni wierzchołek pierwotnego sympleksu; pozostałe ściany i degeneracje zachowujemy przez prawa składania i identyczności. Ponieważ sympleks nerwu jest wyznaczony przez kolejne strzałki, otrzymujemy przedłużenie $K^{\triangleright}\to NI$.

Dla odwrotnej implikacji oraz ostatniej tezy używamy tego samego kryterium filtrowania. Diagram pusty daje obiekt. Dyskretny zbiór symplicjalny mniej niż $\kappa$ wierzchołków daje wspólny cel. Dla mniej niż $\kappa$ par równoległych morfizmów w $h\mathcal I$ o wspólnym celu wybieramy reprezentujące je krawędzie. Tworzymy zbiór symplicjalny z jednym wierzchołkiem odpowiadającym wspólnemu celowi i osobnymi wierzchołkami dla dziedzin par; wstawiamy dwie równoległe krawędzie dla każdej pary. Nie dodajemy sympleksów niezdegenerowanych wyższych wymiarów. Regularność i nieskończoność $\kappa$ zapewniają, że taki zbiór jest $\kappa$-mały. Stożek nad nim daje morfizm ze wspólnego celu, a dwa trójkąty stożka dla każdej pary dają równość złożeń w $h\mathcal I$. Kryterium daje filtrowanie $h\mathcal I$. Dla $\mathcal I=NI$ używamy $hNI\cong I$.
\end{proof}

\begin{przyklad}[Dlaczego kategoria homotopii nie wystarcza]\label{cat:ex:infty-filtered-homotopy-insufficient}
Niech $K=\operatorname{Sing}(S^2)$ będzie kompleksem Kana. Jego kategoria homotopii jest równoważna kategorii końcowej, więc jest filtrowana. Sam $K$ nie jest filtrowany: wybieramy odwzorowanie $\partial\Delta^3\to K$ odpowiadające ciągłemu odwzorowaniu $|\partial\Delta^3|\to S^2$ stopnia jeden. Gdyby istniało przedłużenie na stożek górny, po realizacji i zastosowaniu ewaluacji $|\operatorname{Sing}(S^2)|\to S^2$ otrzymalibyśmy przedłużenie tego odwzorowania na stożek nad sferą. Taki stożek jest ściągalny, więc ograniczenie byłoby homotopijne z odwzorowaniem stałym, którego stopień jest zerowy. Jest to sprzeczność. Filtrowanie jest zatem warunkiem rzeczywiście wyższym.
\end{przyklad}

\subsection{Zwartość względem kogranic filtrowanych}

\begin{definicja}[Obiekt $\kappa$-zwarty]\label{cat:def:infty-kappa-compact}
Niech $\mathcal C$ ma małe $\kappa$-filtrowane kogranice. Obiekt $A$ nazywamy $\kappa$-zwartym, jeśli dla każdego diagramu $D:\mathcal I\to\mathcal C$ o małej $\kappa$-filtrowanej dziedzinie kanoniczne odwzorowanie
\begin{equation}\label{cat:eq:infty-kappa-compact}
 \varinjlim_{i\in\mathcal I}\operatorname{Map}_{\mathcal C}(A,D_i)
 \longrightarrow
 \operatorname{Map}_{\mathcal C}\left(A,\varinjlim_{i\in\mathcal I}D_i\right)
\end{equation}
jest równoważnością przestrzeni. Dla $\kappa=\aleph_0$ mówimy o obiekcie zwartym. Przez $\mathcal C^\kappa$ oznaczamy pełną podkategorię obiektów $\kappa$-zwartych.
\end{definicja}

Kanoniczne odwzorowanie pochodzi ze stożka kogranicznego diagramu $D$: złożenie z jego strzałkami daje zgodną rodzinę odwzorowań przestrzeni, a własność uniwersalna kogranicy w $\mathcal S$ daje~\eqref{cat:eq:infty-kappa-compact}. Jest to odpowiednik~\eqref{cat:eq:accessible-hom-colim} z teorii zwykłych kategorii. Zastosowanie $\pi_0$ daje warunek faktoryzacji klas morfizmów, ale pełna definicja kontroluje również wszystkie wyższe homotopie.

\begin{twierdzenie}[Kogranice filtrowane przestrzeni; wynik zewnętrzny]\label{cat:thm:infty-filtered-space-limits}
W $\mathcal S$ małe $\kappa$-filtrowane kogranice komutują z $\kappa$-małymi granicami (indeksowanymi zbiorami symplicjalnymi o mniej niż $\kappa$ niezdegenerowanych sympleksów). W szczególności zwykłe filtrowane kogranice przestrzeni komutują ze skończonymi granicami.
\end{twierdzenie}

Źródło: J.~Lurie, \emph{Higher Topos Theory}, \S5.3.3, \url{https://www.math.ias.edu/~lurie/papers/HTT.pdf}. Przyjmujemy ten wynik jako twierdzenie zewnętrzne. Nie zastępujemy go twierdzeniem o kogranicach zbiorów: granice przestrzeni obejmują na przykład homotopijne iloczyny włókniste, których zbiór składowych nie musi być zwykłym iloczynem włóknistym zbiorów składowych.

\begin{stwierdzenie}[Domknięcie obiektów zwartych]\label{cat:prop:infty-kappa-compact-closure}
Jeżeli w $\mathcal C$ istnieją $\kappa$-małe kogranice, to $\mathcal C^\kappa$ jest na nie domknięta. Jest także domknięta na równoważności i retrakty w $h\mathcal C$.
\end{stwierdzenie}
\begin{proof}
Niech $A=\varinjlim_{j\in J}A_j$, gdzie $J$ jest $\kappa$-mały i wszystkie $A_j$ są $\kappa$-zwarte. Dla $\kappa$-filtrowanego małego diagramu $D:\mathcal I\to\mathcal C$ mamy równoważności
\[
\begin{aligned}
 \operatorname{Map}\left(A,\varinjlim_iD_i\right)
 &\simeq\varprojlim_{j\in J^{\mathrm{op}}}\operatorname{Map}\left(A_j,\varinjlim_iD_i\right)\\
 &\simeq\varprojlim_{j\in J^{\mathrm{op}}}\varinjlim_i\operatorname{Map}(A_j,D_i)\\
 &\simeq\varinjlim_i\varprojlim_{j\in J^{\mathrm{op}}}\operatorname{Map}(A_j,D_i)\\
 &\simeq\varinjlim_i\operatorname{Map}(A,D_i).
\end{aligned}
\]
Pierwsza i ostatnia równoważność wynikają z własności uniwersalnej kogranicy $A$, druga ze zwartości $A_j$, a trzecia z twierdzenia~\ref{cat:thm:infty-filtered-space-limits}. Wszystkie są naturalne w diagramach, więc ich złożenie jest odwrotnością kanonicznego odwzorowania~\eqref{cat:eq:infty-kappa-compact}. W szczególności obiekt początkowy jest zwarty, jeśli istnieje: jest kogranicą pustego diagramu.

Równoważność obiektów indukuje naturalną równoważność reprezentowanych funktorów, więc zachowuje zwartość. Jeśli $A$ jest retraktem $B$ w $h\mathcal C$, to przez kontrawariantną funktorialność $\operatorname{Map}(-,-)$ otrzymujemy retrakt funktora $\operatorname{Map}(A,-)$ funktora $\operatorname{Map}(B,-)$ w kategorii homotopii kategorii funktorów. Naturalność porównania~\eqref{cat:eq:infty-kappa-compact} daje retrakt odpowiedniego porównania dla $B$ w kategorii homotopii kategorii strzałek przestrzeni. Retrakt izomorfizmu jest izomorfizmem, jak wykazano w dowodzie stwierdzenia~\ref{cat:prop:infty-localization-saturation}. Jeśli $B$ jest $\kappa$-zwarty, to zatem także $A$.
\end{proof}

\subsection{Dostępność i prezentowalność}

\begin{definicja}[Dostępna i prezentowalna $\infty$-kategoria]\label{cat:def:infty-presentable}
$\infty$-kategorię $\mathcal C$ nazywamy $\kappa$-dostępną, jeśli ma małe $\kappa$-filtrowane kogranice i istnieje mała pełna podkategoria $\mathcal A$ złożona z obiektów $\kappa$-zwartych taka, że każdy obiekt $\mathcal C$ jest równoważny kogranicy pewnego małego $\kappa$-filtrowanego diagramu w $\mathcal A$. Nazywamy ją dostępną, jeśli istnieje taka regularna liczba kardynalna $\kappa$. Dostępna $\infty$-kategoria mająca wszystkie małe kogranice jest prezentowalna.
\end{definicja}

Małość rodziny generatorów nie oznacza małości całej kategorii. Na przykład w zwykłej kategorii zbiorów dowolny zbiór jest sumą swoich skończonych podzbiorów, podczas gdy nie istnieje zbiór wszystkich zbiorów w tym samym uniwersum. W wyższej teorii analogiczną rolę pełnią odpowiednie kogranice przestrzeni i obiekty $\kappa$-zwarte.

\begin{stwierdzenie}[Zwarte obiekty jako retrakty generatorów]\label{cat:prop:infty-accessible-compact-retract}
Jeżeli $\mathcal C$ jest $\kappa$-dostępna z podkategorią $\mathcal A$ jak w definicji~\ref{cat:def:infty-presentable}, to każdy obiekt $A\in\mathcal C^\kappa$ jest retraktem w $h\mathcal C$ pewnego obiektu $\mathcal A$.
\end{stwierdzenie}
\begin{proof}
Piszemy $A\simeq\varinjlim_{i\in\mathcal I}A_i$, gdzie $A_i\in\mathcal A$. Oznaczamy strzałki stożka przez $u_i:A_i\to A$. Przez $\kappa$-zwartość $A$ odwzorowanie~\eqref{cat:eq:infty-kappa-compact} jest równoważnością. Funktor $\pi_0:\mathcal S\to\mathrm{Set}$ zachowuje kogranice, gdyż jest lewym sprzężonym do włączenia zbiorów jako przestrzeni dyskretnych. Otrzymujemy bijekcję
\[
 \varinjlim_{i\in h\mathcal I}\Hom_{h\mathcal C}(A,A_i)
 \cong\Hom_{h\mathcal C}(A,A).
\]
Przejście do kategorii homotopii indeksu jest uzasadnione tym, że zgodność rodzin funkcji do zbioru dyskretnego zależy tylko od klas krawędzi (por. dowód wniosku~\ref{cat:cor:infty-verdier-mapping-groups}). Każdy element zwykłej kogranicy zbiorów ma reprezentanta w jednym z jej składników, ponieważ kogranica jest ilorazem ich rozłącznej sumy. W szczególności $\id_A$ jest reprezentowana przez klasę $v:A\to A_i$. Zatem $[u_i][v]=\id_A$ w $h\mathcal C$, co daje żądany retrakt.
\end{proof}

\begin{uwaga}[Dalszy zakres]\label{cat:rem:infty-presentable-scope}
Podstawowe twierdzenia o prezentowalności (opis przez dostępne lokalizacje kategorii presnopów przestrzeni, zupełność i kryterium istnienia sprzężeń) rozwiniemy w następnych podrozdziałach. Następnie rozpatrzymy lokalizacje Bousfielda, nasycenia klas morfizmów, obiekty lokalne oraz przypadek stabilny. Ogólna teoria nie wymaga istnienia konkretnej struktury modelowej. Pojęcia z rozdziału~\ref{cat:sec:derived-localization-morita} będą używane przez ich własności uniwersalne.
\end{uwaga}
